% CHAPTER 4: SYSTEM REQUIREMENTS AND SPECIFICATION
\chapter{System Requirements and Specification}

\section{Hardware Requirements}

\subsection{Microcontroller -- ESP32}
The ESP32 will serve as the central processing and communication unit of the RespiraTech system. It is based on a dual-core Xtensa LX6 processor running at 240 MHz with 520 KB of SRAM, 4 MB of flash memory, and integrated Wi-Fi (IEEE 802.11 b/g/n) and Bluetooth 4.2 connectivity. The ESP32 features 18 channels of 12-bit SAR (Successive Approximation Register) analog-to-digital converters (ADCs), which will be used to read the analog voltage outputs of the MQ-series gas sensors. The ADC provides a resolution of 4096 levels across a 0--3.3V input range, corresponding to a voltage resolution of approximately 0.8 mV per step. The module also includes multiple GPIO pins for digital I/O, I2C communication (for the OLED display), and SPI interfaces. The integrated Wi-Fi stack will enable direct HTTP/HTTPS communication with cloud platforms without requiring external networking hardware.

\subsection{Gas Sensor -- MQ-135 (Air Quality Sensor)}
The MQ-135 is a metal-oxide semiconductor (MOS) gas sensor designed for air quality monitoring. It is sensitive to ammonia (NH$_3$), nitrogen oxides (NO$_x$), benzene (C$_6$H$_6$), carbon dioxide (CO$_2$), alcohol vapours, and smoke. The sensor operates on a heating element principle: a tin dioxide (SnO$_2$) semiconductor film changes its electrical resistance when exposed to target gases. The sensor requires a 5V heater supply and outputs an analog voltage proportional to the gas concentration through a voltage divider circuit with a load resistor ($R_L$). For breath analysis, the MQ-135 will detect broad-spectrum VOC concentrations in the exhaled sample, providing data on air quality markers that correlate with respiratory conditions and metabolic disorders. The sensor's preheat time is approximately 24--48 hours for initial calibration, with a response time of less than 10 seconds during operation.

\subsection{Gas Sensor -- MQ-3 (Alcohol/Acetone Sensor)}
The MQ-3 is a semiconductor gas sensor specifically designed for alcohol and acetone vapour detection. It uses a SnO$_2$ semiconductor element that exhibits high sensitivity to ethanol (C$_2$H$_5$OH) and acetone ((CH$_3$)$_2$CO) with a detection range of 0.04--4 mg/L for alcohol. The sensor will be critical for diabetes screening, as elevated acetone levels in breath (above 1.8 ppm) are a well-established biomarker for diabetic ketoacidosis. Like the MQ-135, it requires a 5V heater supply and provides an analog output voltage. The complementary sensitivities of the MQ-135 (broad-spectrum) and MQ-3 (acetone-specific) will create a sensor array capable of capturing a broader VOC profile than either sensor alone.

\subsection{Temperature and Humidity Sensor -- DHT11/DHT22}
The DHT11 (or upgraded DHT22) sensor will measure ambient temperature and relative humidity to establish environmental baseline conditions. Gas sensor readings are significantly affected by temperature (affecting molecular kinetic energy and sensor surface reactivity) and humidity (affecting water vapour interference with target gas adsorption). The DHT11 provides temperature measurements from 0--50°C with $\pm$2°C accuracy and humidity measurements from 20--80\% RH with $\pm$5\% accuracy. The upgraded DHT22 offers improved ranges (--40 to 80°C, 0--100\% RH) and better accuracy ($\pm$0.5°C, $\pm$2\% RH). The sensor communicates via a single-wire digital protocol, requiring only one GPIO pin on the ESP32. Environmental data will be used to normalize gas sensor readings and compensate for ambient conditions, improving the reliability of VOC measurements.

\subsection{Display Module -- SSD1306 OLED (0.96 inch, 128$\times$64, I2C)}
The SSD1306 OLED display will provide local visual feedback to the user during the breath analysis process. It features a 128$\times$64 pixel resolution on a 0.96-inch diagonal screen and communicates with the ESP32 via the I2C protocol at device address 0x3C. The display will show real-time sensor readings during the breath capture phase, a progress indicator during cloud processing, and a summary of the AI prediction results including disease risk percentages. The OLED technology provides high contrast, wide viewing angles, and low power consumption (approximately 20 mA at 3.3V), making it suitable for battery-powered portable operation.

\subsection{Breath Collection Chamber}
A dedicated breath collection chamber will be designed to provide a controlled environment for the breath sample to interact with the gas sensors. The chamber will be a small enclosed volume (approximately 50--100 mL) with a mouthpiece inlet for the user to exhale into and a ventilation outlet with a one-way valve to prevent ambient air from entering. The MQ-135 and MQ-3 sensors will be mounted inside the chamber with their sensing elements exposed to the enclosed air volume. The chamber will be designed to be easily detachable and cleanable to maintain hygiene between uses.

\subsection{Power Supply -- USB/Battery}
The system will be powered via a micro USB connection for desktop/stationary use, or a 3.7V lithium polymer battery (1000--2000 mAh) with a TP4056 USB charging module for portable operation. The ESP32 operates at 3.3V (supplied via its onboard voltage regulator from the USB 5V input), while the MQ-series gas sensors require 5V for their heating elements. A boost converter or direct USB 5V supply will provide the heater voltage. Two 100nF ceramic decoupling capacitors will be placed near the ESP32 power input and sensor power rails to filter high-frequency noise.

\section{Software Requirements}

\subsection{Programming Languages}
\begin{itemize}[leftmargin=1.5cm]
\item \textbf{C/C++ (Arduino Framework):} Will be used for ESP32 firmware development within the Arduino IDE 2.x environment. The firmware will handle ADC sampling of gas sensors, DHT sensor reading, I2C OLED display communication, Wi-Fi connectivity, HTTP/HTTPS client operations for cloud data transmission, and JSON payload construction.

\item \textbf{Python 3.10+:} Will serve as the primary language for the machine learning pipeline, including data preprocessing, feature engineering, model training (using scikit-learn and TensorFlow/Keras), model evaluation, and cloud function development for inference endpoints.

\item \textbf{JavaScript/HTML/CSS:} Will be used for the mobile/web application frontend that displays health risk predictions, historical data visualizations (using Chart.js or D3.js), and alert notifications.
\end{itemize}

\subsection{Cloud Platforms and Libraries}
\begin{itemize}[leftmargin=1.5cm]
\item \textbf{ThingSpeak:} An IoT analytics cloud platform that will receive sensor data from the ESP32 via HTTP API calls, store it in channels with timestamped fields, and provide MATLAB-based data analysis and visualization capabilities. ThingSpeak supports up to 8 data fields per channel and allows data updates every 15 seconds on the free tier.

\item \textbf{Firebase (alternative):} Google's mobile and web application platform that will provide Realtime Database or Firestore for sensor data storage, Cloud Functions for serverless ML inference, and Firebase Hosting for the web application frontend.

\item \textbf{scikit-learn:} A Python machine learning library that will provide implementations of Random Forest, SVM, and other classification algorithms for disease prediction model training and evaluation.

\item \textbf{TensorFlow/Keras:} A deep learning framework that will be used to develop and train neural network classifiers for more complex pattern recognition in the sensor data, potentially achieving higher accuracy than traditional ML algorithms.

\item \textbf{Flask/FastAPI:} A Python web framework that will host the ML inference API endpoint if a custom cloud server is used instead of ThingSpeak/Firebase.
\end{itemize}

\subsection{Communication Protocols}
\begin{itemize}[leftmargin=1.5cm]
\item \textbf{Wi-Fi (IEEE 802.11 b/g/n):} The primary wireless communication protocol for transmitting sensor data from the ESP32 to the cloud platform over the local area network and internet.

\item \textbf{HTTP/HTTPS:} The application-layer protocol for RESTful API communication between the ESP32 and cloud platforms. HTTPS (TLS 1.2+) will be used for secure data transmission to protect sensitive health information.

\item \textbf{I2C (Inter-Integrated Circuit):} The two-wire serial protocol connecting the ESP32 to the SSD1306 OLED display at device address 0x3C with 400 kHz clock speed.

\item \textbf{Single-Wire Protocol:} The proprietary communication protocol used by the DHT11/DHT22 sensor for transmitting temperature and humidity data to the ESP32.

\item \textbf{MQTT (optional):} A lightweight publish-subscribe messaging protocol that may be used as an alternative to HTTP for real-time data streaming between the ESP32 and cloud broker, offering lower overhead and persistent connections.
\end{itemize}

\section{Power Requirements}
The system will operate on USB 5V power or a 3.7V LiPo battery with the following estimated power budget:

\begin{table}[h]
\centering
\begin{tabular}{|l|c|c|}
\hline
\textbf{Component} & \textbf{Current Draw} & \textbf{Voltage} \\
\hline
ESP32 (active with Wi-Fi) & 160--240 mA & 3.3V \\
MQ-135 sensor (heater active) & 150 mA & 5.0V \\
MQ-3 sensor (heater active) & 150 mA & 5.0V \\
DHT11/DHT22 sensor & 2.5 mA & 3.3V \\
SSD1306 OLED Display & 20 mA & 3.3V \\
\hline
\textbf{Total Peak Consumption} & \textbf{$\approx$573 mA} & \textbf{mixed} \\
\hline
\end{tabular}
\caption{Estimated Power Consumption Budget}
\end{table}

The gas sensor heaters are the dominant power consumers. With a 2000 mAh battery, the system is expected to provide approximately 3--3.5 hours of continuous operation. In practical intermittent usage scenarios (where sensors are preheated once and individual breath samples are taken periodically), battery life is expected to extend to approximately 4--6 hours.
